\documentclass[11pt,a4paper]{coverletter}

\senderblock
  {Radheem Bin Razi}
  {Distributed Systems \& Cloud-Native Engineer}
  {Ilmenau, Germany \textbar\ \href{mailto:sheikh.radheem@gmail.com}{sheikh.radheem@gmail.com} \textbar\ \href{https://radheem.github.io/my-cv}{portfolio}}

\begin{document}

%% ===== ENGLISH =====
\recipient{1KOMMA5°}{Hiring Team}{}

\opening{Dear Hiring Team,}

\vspace{8pt}\noindent\textbf{1. Why 1KOMMA5°?}\par\vspace{3pt}

A reliable device-to-cloud pipeline keeps the lights on long before a dashboard ever flashes a warning. Architecting that quiet resilience for physical infrastructure is exactly where I want to focus my career. I am applying for the IoT Cloud Platform Engineer role because 1KOMMA5° is building the messaging backbone that will connect and manage tens of thousands of residential energy devices across Germany. The chance to shape a Kubernetes-native platform that routes MQTT telemetry while feeding real-time fleet insights aligns perfectly with my background in distributed systems.

\vspace{8pt}\noindent\textbf{2. Why me?}\par\vspace{3pt}

In my recent distributed systems work, I replaced a heavy sidecar architecture with native NATS JetStream for high-throughput messaging. This change simplified routing and reduced operational overhead while preserving delivery guarantees. I deployed those workloads on Kubernetes using GitOps tooling and Cilium network policies. I implemented automated health checks that caught configuration drift before it impacted production traffic. I am comfortable hardening these systems with tracing and setting up alerting thresholds. I also write Go services that scale predictably under load.

That same reliability mindset drives how I build data pipelines. I have orchestrated multi-step ingestion workflows with declarative routing and automated quality gates. This ensures that telemetry and metadata land in PostgreSQL and document stores without silent failures. I enjoy designing the exact seam between device gateways and cloud services. Clean message schemas and observability pipelines turn raw telemetry into actionable fleet insights. I am ready to bring that operational rigor to your GKE environment and Eclipse Hono stack.

\vspace{8pt}\noindent\textbf{3. Why now?}\par\vspace{3pt}

I am looking to settle into a role where I can own the full lifecycle of a critical cloud platform. 1KOMMA5° offers exactly that scope. My five years of building and operating production backends have prepared me to step in immediately. I will contribute directly to your gateway fleet and observability pipelines. I am available for full-time work starting now. I am based in Germany with the ability to relocate within the country within two to three weeks. I hold Blue Card eligibility. Your team can proceed with the contract without any sponsorship overhead.

\closing{Sincerely,}{Radheem Bin Razi}

\clearpage
\selectlanguage{ngerman}

%% ===== DEUTSCH =====
\recipient{1KOMMA5°}{Personalabteilung}{}

\opening{Sehr geehrtes Hiring-Team,}

\vspace{8pt}\noindent\textbf{1. Warum 1KOMMA5°?}\par\vspace{3pt}

Eine zuverlässige Device-to-Cloud-Pipeline sorgt dafür, dass die Systeme lange vor der ersten Warnmeldung im Dashboard noch stabil laufen. Die Architektur einer solchen stillen Resilienz für physische Infrastruktur ist genau der Bereich, in dem ich meine Karriere fokussieren möchte. Ich bewerbe mich für die Position als IoT Cloud Platform Engineer, da 1KOMMA5° das Messaging-Backbone aufbaut, das zehntausende von Energiegeräten in ganz Germany vernetzen und verwalten wird. Die Chance, eine Kubernetes-native Plattform zu gestalten, die MQTT-Telemetrie routet und gleichzeitig Echtzeit-Einblicke in die Geräteflotte liefert, entspricht perfekt meinem Hintergrund in verteilten Systemen.

\vspace{8pt}\noindent\textbf{2. Warum ich?}\par\vspace{3pt}

In meinen aktuellen Projekten mit verteilten Systemen habe ich eine schwere Sidecar-Architektur durch natives NATS JetStream für High-Throughput-Messaging ersetzt. Diese Änderung vereinfachte das Routing und reduzierte den Betriebsaufwand, während die Delivery-Garantien erhalten blieben. Ich habe diese Workloads auf Kubernetes mit GitOps-Tooling und Cilium network policies deployed. Ich habe automatisierte Health Checks implementiert, die Konfigurationsdrift erkannt haben, bevor sie den Production-Traffic beeinträchtigt hat. Ich bin erfahren darin, diese Systeme durch Tracing zu härten und Alerting-Schwellenwerte einzurichten. Zudem schreibe ich Go-Services, die unter Last vorhersagbar skalieren.

Dieselbe Zuverlässigkeitslogik treibt auch meinen Ansatz beim Aufbau von Datenpipelines an. Ich habe mehrstufige Ingestion-Workflows mit deklarativem Routing und automatisierten Quality Gates orchestriert. Dies stellt sicher, dass Telemetrie und Metadaten ohne stille Fehler in PostgreSQL und document stores landen. Ich gestalte gerne die genaue Schnittstelle zwischen device gateways und cloud services. Saubere message schemas und observability pipelines verwandeln rohe Telemetrie in umsetzbare Einblicke in die Geräteflotte. Ich bin bereit, diese operative Sorgfalt in Ihre GKE-Umgebung und den Eclipse Hono stack einzubringen.

\vspace{8pt}\noindent\textbf{3. Warum jetzt?}\par\vspace{3pt}

Ich suche eine Stelle, in der ich den kompletten Lebenszyklus einer kritischen Cloud-Plattform verantworten kann. 1KOMMA5° bietet genau diesen Verantwortungsbereich. Meine fünf Jahre Erfahrung im Aufbau und Betrieb von production backends haben mich darauf vorbereitet, mich sofort einzubringen. Ich werde mich direkt an Ihrer gateway fleet und den observability pipelines beteiligen. Ich stehe ab sofort für eine Vollzeitbeschäftigung zur Verfügung. Ich lebe in Germany und kann mich innerhalb von zwei bis drei Wochen innerhalb des Landes umziehen. Ich erfülle die Voraussetzungen für eine Blue Card. Ihr Team kann den Vertrag ohne zusätzlichen Sponsorship-Aufwand fortführen.

\closing{Mit freundlichen Grüßen,}{Radheem Bin Razi}

\end{document}
